%!TEX root = ../thesis.tex
\mychapter{System Design}{System Design}{}
\label{chap:design}

\mysection{User and functional requirements}{User and Functional Requirements}
\label{sec:requirements}

\mysection{System architecture}{System Architecture}
\label{sec:architecture}

\mysubsection{Off-board inference and the thin-client split}{Off-board Inference}
\label{ssec:split}

\mysubsection{Two-rate control: waypoint intents versus the reactive layer}{Two-rate Control}
\label{ssec:rates}

\mysection{The capability vocabulary}{The Capability Vocabulary}
\label{sec:vocabulary}

\mysection{Clarification loop design}{Clarification Loop Design}
\label{sec:clarification_design}

\mysection{Navigation as instrumentation}{Navigation as Instrumentation}
\label{sec:navigation}

\mysubsection{The line follower as the reactive layer}{The Line Follower}
\label{ssec:line_follower}

Section~\ref{ssec:rates} showed that the Mac-hosted models cannot sit in a
steering loop. A VLM call takes about 15\,s, and keeping a mecanum chassis on a
line needs control at 10 to 30\,Hz. I therefore left steering to the TurboPi's
stock \texttt{line\_follow} ROS2 node, which closes the loop on the Pi using the
chassis-mounted 4-channel IR array. I used the IR array instead of the gimbal
camera so the camera stays free for the VLM to look wherever the mission needs,
regardless of where the chassis is steering (\texttt{rover\_pi.py},
\texttt{aim()}).

\texttt{PiRoverController} does not reimplement steering. It only supervises the
node on three things the node has no logic for itself:
\begin{itemize}
  \item Arrival: the node's own \texttt{/line\_follow/crossroads\_stop} topic
    decides when the end of the line is reached, not a model's judgement.
  \item Obstacles: the ultrasonic sensor is polled independently of the line
    sensor. Two consecutive readings below the stop threshold pause the follower
    until the path clears, and then it resumes.
  \item Giving up: a wall-clock timeout. The stock node has no lost-line exit,
    so without one the rover would drive straight on indefinitely.
\end{itemize}

I also use a modified variant, \texttt{line\_follow\_corner.py}. When it takes a
90\textdegree{} bend instead of parking at a crossbar, it publishes a
\texttt{/line\_follow/corner} event. This lets a route be given as an ordered
string of turns (e.g.\ \texttt{"L"}, \texttt{"R"}) and not only the
straight-track default. It adds one constraint: a planned pause, such as a
mid-drive look, must not land while the follower is mid-pivot. A \texttt{STOP}
sent during the node's pre-turn nudge cancels the turn instead of the forward
motion. \texttt{corner\_guard\_active()} reports a short window after each bend,
and callers must defer their pause until it ends.

Halting the follower is where a naive design fails without any error. While
running, the node republishes \texttt{/cmd\_vel} at 100\,Hz, so a single
zero-velocity \texttt{Twist} is overwritten within 10\,ms and the rover keeps
driving. \texttt{halt()} therefore sends an explicit \texttt{STOP} service call
to the node \emph{before} the zero \texttt{Twist}. It repeats both 0.7\,s later,
in case the first \texttt{STOP} landed inside the node's pre-turn nudge window
and was swallowed. This is the most safety-critical part of the hardware
integration.

\mysection{Summary of design decisions}{Summary of Design Decisions}
\label{sec:design_decisions}
